\documentclass[conference]{IEEEtran}
\IEEEoverridecommandlockouts
% The preceding line is only needed to identify funding in the first footnote. If that is unneeded, please comment it out.
\usepackage{cite}
\usepackage{amsmath,amssymb,amsfonts}
\usepackage{algorithmic}
\usepackage{graphicx}
\usepackage{textcomp}
\usepackage{xcolor}
\usepackage[brazilian]{babel}
\usepackage[utf8]{inputenc}
\def\BibTeX{{\rm B\kern-.05em{\sc i\kern-.025em b}\kern-.08em
    T\kern-.1667em\lower.7ex\hbox{E}\kern-.125emX}}

% Texto em vermelho para marcar o que ainda falta escrever (remover ao final)
\newcommand{\pendente}[1]{{\color{red}[#1]}}

\begin{document}


\title{Trabalho Computacional AV2: Redes Neurais Artificiais \\
{\footnotesize T296 - Inteligência Artificial Computacional}}

\author{
\IEEEauthorblockN{Gabriel de Sousa Nobre}
\IEEEauthorblockA{
\textit{Graduação em Ciência da Computação} \\
\textit{Universidade de Fortaleza (UNIFOR)} \\
Fortaleza, Ceará, Brasil \\
gabrielnobre.dev@gmail.com
}
\and
\IEEEauthorblockN{Victor Lins Gurgel do Amaral}
\IEEEauthorblockA{
\textit{Graduação em Ciência da Computação} \\
\textit{Universidade de Fortaleza (UNIFOR)} \\
Fortaleza, Ceará, Brasil \\
vlinsgurgel@gmail.com
}
\and
\IEEEauthorblockN{Lorenzo Barros Calheiros Pinheiro}
\IEEEauthorblockA{
\textit{Graduação em Ciência da Computação} \\
\textit{Universidade de Fortaleza (UNIFOR)} \\
Fortaleza, Ceará, Brasil \\
barroslorenzopro@gmail.com
}}

\maketitle

\begin{abstract}
Este trabalho apresenta a implementação e a avaliação de redes neurais
artificiais em duas etapas.
\pendente{Primeira etapa: resumo da classificação do conjunto spiral\_d com
Perceptron Simples, ADALINE e MLP.}
Na segunda etapa, os modelos Perceptron Simples, ADALINE e Perceptron de
Múltiplas Camadas (MLP) foram aplicados ao reconhecimento facial de 20
pessoas, a partir de 640 imagens redimensionadas para 30 por 30 pixels. Os
modelos foram implementados apenas com a biblioteca NumPy e validados por
simulações de Monte Carlo com 10 rodadas e partição 80/20. A MLP obteve a
maior acurácia média (97,42\%), seguida de perto pelo Perceptron Simples
(96,80\%), enquanto o ADALINE obteve 92,50\% e a maior variação entre as
rodadas. Os erros se concentraram em poucas pessoas e imagens, e os três
modelos erraram repetidamente algumas das mesmas imagens.
\end{abstract}

\begin{IEEEkeywords}
redes neurais artificiais, perceptron simples, ADALINE, perceptron de
múltiplas camadas, reconhecimento facial, simulação de Monte Carlo
\end{IEEEkeywords}

\section{Introdução}

As redes neurais artificiais (RNA) são modelos computacionais inspirados no
sistema nervoso. Elas são formadas por neurônios artificiais ligados por
pesos sinápticos, que são ajustados a partir de exemplos durante o
treinamento~\cite{b1}. Neste trabalho foram implementados os modelos
Perceptron Simples~\cite{b2}, ADALINE~\cite{b3} e Perceptron de Múltiplas
Camadas (MLP), treinado pelo algoritmo de retropropagação do
erro~\cite{b4}. Foram usadas apenas
as bibliotecas NumPy, Matplotlib, Seaborn e OpenCV.

O trabalho foi dividido em duas etapas. Na primeira, os modelos foram
aplicados a um problema de classificação bidimensional com duas classes,
que serve para validar as implementações. Na segunda, os mesmos modelos foram
aplicados ao reconhecimento facial de 20
pessoas, um problema com 20 classes e centenas de atributos. As
Seções~\ref{sec:met1} e~\ref{sec:res1} tratam da primeira etapa, as
Seções~\ref{sec:met2} e~\ref{sec:res2} tratam da segunda, e a
Seção~\ref{sec:conclusao} apresenta as conclusões.

% =====================================================================
% PRIMEIRA ETAPA (a ser escrita pela equipe)
% =====================================================================

\section{Classificação Bidimensional: Metodologia}\label{sec:met1}

\subsection{Organização dos dados}
\pendente{Organização do conjunto spiral\_d.csv (1400 amostras, duas
classes rotuladas na terceira coluna) e normalização.}

\subsection{Visualização inicial}
\pendente{Gráfico de espalhamento dos dados.}

\subsection{Modelos implementados}
\pendente{Perceptron Simples, ADALINE e MLP.}

\subsection{Definição dos hiperparâmetros}
\pendente{Como os hiperparâmetros de cada modelo foram escolhidos.}

\subsection{\textit{Underfitting} e \textit{overfitting}}
\pendente{Topologias subdimensionadas e superdimensionadas da rede MLP.}

\subsection{Validação e métricas}
\pendente{Amostragem aleatória com $R = 500$ rodadas (80\% treino e 20\%
teste); acurácia, sensibilidade, especificidade, precisão e F1-score.}

\section{Classificação Bidimensional: Resultados e Discussão}\label{sec:res1}

\subsection{Desempenho dos modelos}
\pendente{Uma tabela por métrica (média, desvio-padrão, maior e menor
valor), boxplot ou violin plot e discussão.}

% Modelo de tabela (uma para cada métrica):
% \begin{table}[htbp]
% \caption{Estatísticas da acurácia ao longo das 500 rodadas}
% \begin{center}
% \begin{tabular}{|l|c|c|c|c|}
% \hline
% \textbf{Modelos} & \textbf{Média} & \textbf{Desvio-Padrão} & \textbf{Maior} & \textbf{Menor} \\
% \hline
% Perceptron Simples & & & & \\
% \hline
% ADALINE & & & & \\
% \hline
% MLP & & & & \\
% \hline
% \end{tabular}
% \label{tab:acc1}
% \end{center}
% \end{table}

\subsection{Matrizes de confusão e curvas de aprendizado}
\pendente{Rodadas de maior e menor valor de cada métrica.}

\subsection{\textit{Underfitting} e \textit{overfitting}}
\pendente{Desempenho das topologias com underfitting e overfitting da rede
MLP.}

% =====================================================================
% SEGUNDA ETAPA: RECONHECIMENTO FACIAL
% =====================================================================

\section{Reconhecimento Facial: Metodologia}\label{sec:met2}

\subsection{Organização dos dados}

O conjunto RecFac possui 640 imagens em tons de cinza de 20 pessoas, com 32
imagens por pessoa e dimensão de $120 \times 128$ pixels. As imagens de cada
pessoa variam a direção do rosto (frente, esquerda, direita e para cima), a
expressão facial (neutra, feliz, triste e com raiva) e o uso de óculos
escuros. Esse conjunto corresponde ao CMU Face Images~\cite{b5}.

Cada imagem foi lida com a biblioteca OpenCV e redimensionada para
$30 \times 30$ pixels (Fig.~\ref{fig:redim}), a menor das opções do
enunciado. Assim, cada imagem passa a ter $p = 900$ atributos, em vez de
15.360, o que reduz bastante o custo do treinamento, e os rostos continuam
distinguíveis. Cada imagem foi transformada em um vetor coluna, e os
vetores foram reunidos na matriz $\mathbf{X} \in \mathbb{R}^{p \times N}$,
com $N = 640$.

\begin{figure}[htbp]
\centerline{\includegraphics[width=\columnwidth]{exemplo_redimensionamento.png}}
\caption{Exemplos de imagens do conjunto RecFac no tamanho original (acima)
e redimensionadas para $30 \times 30$ pixels (abaixo).}
\label{fig:redim}
\end{figure}

Os pixels foram normalizados para o intervalo $[-1, +1]$ pelo método
min-max, indicado nas notas de aula para redes que usam a tangente
hiperbólica~\cite{b1}:
\begin{equation}
x^* = 2 \cdot \frac{x - \min(\mathbf{X})}{\max(\mathbf{X}) - \min(\mathbf{X})} - 1.
\label{eq:minmax}
\end{equation}
Como todos os pixels têm a mesma escala (0 a 255), foram usados o menor e o
maior valor de toda a matriz. Em seguida, foi adicionada uma linha de
valores $-1$ (o bias) no topo da matriz, resultando em
$\mathbf{X} \in \mathbb{R}^{(p+1) \times N} = \mathbb{R}^{901 \times 640}$.

Os rótulos foram codificados com \textit{one-hot encoding} bipolar: cada
pessoa é representada por um vetor com $+1$ na posição da sua classe e $-1$
nas demais. A matriz de rótulos é, então,
$\mathbf{Y} \in \mathbb{R}^{C \times N} = \mathbb{R}^{20 \times 640}$, com as
pessoas ordenadas pelo nome da pasta (an2i, at33, boland, bpm, ch4f, ...).

\subsection{Modelos implementados}

Os três modelos têm $m = C = 20$ neurônios na camada de saída, um para
cada pessoa. Em todos, os pesos foram iniciados com valores aleatórios no
intervalo $[-0{,}5;\ 0{,}5]$ e o treinamento foi feito amostra por amostra,
com a ordem das amostras embaralhada a cada época. Como mais de um neurônio
de saída (ou nenhum) pode indicar $+1$ ao mesmo tempo, a classe predita é a
do neurônio de saída com o maior valor.

\textbf{Perceptron Simples.} Para cada amostra $x(k)$, calcula-se o
potencial de ativação $\mu(k) = \tilde{W}(k)\, x(k)$, em que $\tilde{W}$ é a
matriz de pesos de ordem $m \times (p+1)$, e a saída $y(k) = g(\mu(k))$ pelo
degrau bipolar, aplicado a cada neurônio. Os pesos são ajustados pela regra
de aprendizagem do Perceptron~\cite{b2}:
\begin{equation}
\tilde{W}(k+1) = \tilde{W}(k) + \eta\, e(k)\, x^T(k), \quad
e(k) = d(k) - y(k),
\label{eq:regra}
\end{equation}
em que $\eta$ é a taxa de aprendizagem e $d(k)$ é a saída desejada. O
Perceptron tradicional só para quando não há erros no conjunto de
treinamento. Como pedido no enunciado, foi acrescentado um número máximo
de épocas: o treinamento para quando uma época termina sem erros ou quando
esse limite é atingido.

\textbf{ADALINE.} Tem a mesma estrutura, mas o erro é calculado com o
potencial de ativação, $e(k) = d(k) - \mu(k)$, e os pesos são ajustados pela
regra Delta (LMS)~\cite{b3}, com a mesma forma de~\eqref{eq:regra}. O
treinamento minimiza o erro quadrático médio (EQM), calculado a cada época
sobre as $N$ amostras de treinamento (512) e as $m$ saídas:
\begin{equation}
EQM = \frac{1}{2Nm} \sum_{k=1}^{N} \sum_{l=1}^{m} \left[ d_l(k) - \mu_l(k) \right]^2.
\label{eq:eqm}
\end{equation}
Essa é a expressão do EQM vista em aula para $m$ saídas, dividida também
por $m$, como no código de referência da disciplina, que faz a média do
erro quadrático sobre as amostras e as saídas~\cite{b1}. O treinamento para
quando o EQM fica abaixo da precisão $pr$ ou quando o número máximo de
épocas é atingido.

\textbf{MLP.} Foi usada uma rede com uma camada oculta de $q$ neurônios e a
camada de saída com $m$ neurônios. Todos os neurônios usam a tangente
hiperbólica, na forma apresentada em aula~\cite{b1}:
\begin{equation}
y = g(\mu) = \frac{1 - e^{-\mu}}{1 + e^{-\mu}}, \qquad
g'(\mu) = \frac{1}{2} \left( 1 - y^2 \right).
\label{eq:ativacao}
\end{equation}
Diferente do ADALINE, o erro é calculado com a saída da rede,
$e(k) = d(k) - y(k)$, e os pesos são ajustados pelo algoritmo de
retropropagação do erro (\textit{backpropagation})~\cite{b4}, seguindo a
implementação apresentada em aula~\cite{b1}. Os critérios de parada são os
mesmos do ADALINE, com o EQM~\eqref{eq:eqm} calculado com $y_l(k)$ no lugar
de $\mu_l(k)$.

\subsection{Definição dos hiperparâmetros}

Os hiperparâmetros foram escolhidos com testes preliminares: cada
configuração foi treinada e testada em três partições 80/20 sorteadas, as
mesmas para todas as configurações. A Tabela~\ref{tab:testes} resume esses
testes e a Tabela~\ref{tab:hiper} mostra os valores escolhidos.

\begin{table}[htbp]
\caption{Testes preliminares dos hiperparâmetros (média de três partições)}
\begin{center}
\begin{tabular}{|l|l|c|c|}
\hline
\textbf{Modelo} & \textbf{Configuração} & \textbf{Acurácia} & \textbf{Épocas} \\
\hline
Perceptron & $\eta = 0{,}1$ & 0,9714 & 10 \\
Simples & $\eta = 0{,}01$ & 0,9688 & 14 \\
 & $\eta = 0{,}001$ & 0,9271 & 46 \\
\hline
ADALINE & $\eta = 0{,}01$ & divergiu & 3 \\
 & $\eta = 0{,}001$ & 0,9349 & 100 \\
 & $\eta = 0{,}0001$ & 0,6849 & 100 \\
\hline
MLP & $q = 20$, $\eta = 0{,}01$ & 0,9635 & 191 \\
 & $q = 50$, $\eta = 0{,}01$ & 0,9792 & 99 \\
 & $q = 134$, $\eta = 0{,}01$ & 0,9740 & 121 \\
 & $q = 50$, $\eta = 0{,}1$ & 0,9688 & 53 \\
 & $q = 50$, $\eta = 0{,}001$ & 0,9714 & 200 \\
\hline
\end{tabular}
\label{tab:testes}
\end{center}
\end{table}

\begin{table}[htbp]
\caption{Hiperparâmetros escolhidos}
\begin{center}
\begin{tabular}{|l|c|c|c|c|}
\hline
\textbf{Modelo} & \textbf{Topologia} & $\boldsymbol{\eta}$ & \textbf{Máx. épocas} & $\boldsymbol{pr}$ \\
\hline
Perceptron Simples & 900--20 & 0,01 & 100 & -- \\
\hline
ADALINE & 900--20 & 0,001 & 100 & $10^{-3}$ \\
\hline
MLP & 900--50--20 & 0,01 & 200 & $10^{-3}$ \\
\hline
\end{tabular}
\label{tab:hiper}
\end{center}
\end{table}

No Perceptron Simples, $\eta = 0{,}1$ e $\eta = 0{,}01$ tiveram praticamente
o mesmo desempenho (diferença de uma imagem em 384), e $\eta = 0{,}001$ foi
pior, provavelmente porque, com passos muito pequenos, os pesos finais
ficam próximos dos pesos iniciais aleatórios. Foi mantido $\eta = 0{,}01$.
Em todos os testes o Perceptron terminou o treinamento sem erros antes de
50 épocas, então o limite de 100 épocas serve apenas como garantia de
parada.

O ADALINE foi o modelo mais sensível à taxa de aprendizagem. Como cada
amostra tem 901 entradas, cada ajuste dos pesos é grande, e com
$\eta = 0{,}01$ o EQM cresceu sem controle até o treinamento divergir. Com
$\eta = 0{,}0001$ o aprendizado foi lento demais. Por isso foi escolhido
$\eta = 0{,}001$. A precisão $pr = 10^{-3}$ é o valor sugerido nas
notas de aula~\cite{b1}, e o limite de 100 épocas foi mantido porque a
maior parte da redução do EQM acontece nas primeiras épocas.

Para a MLP, as notas de aula apresentam regras que dão apenas um valor
inicial para $q$~\cite{b1}: a regra do valor médio, $q = (p+m)/2 = 460$; a
da raiz quadrada, $q = \sqrt{p \cdot m} \approx 134$; e a de Kolmogorov,
$q = 2p + 1 = 1801$, em que $m = 20$ é o número de saídas. Foi testado o
valor da raiz quadrada ($q = 134$) e valores menores, já que os demais
deixariam o treinamento muito mais lento. Com $q = 20$ a rede precisou de
quase todas as 200 épocas e teve a menor acurácia, possivelmente porque a
camada oculta era pequena demais. Com $q = 134$ a acurácia foi parecida com
a de $q = 50$, mas o treinamento foi mais lento. Por isso foi escolhido
$q = 50$. Quanto à taxa de aprendizagem, $\eta = 0{,}1$ convergiu mais
rápido (53 épocas), mas com acurácia menor, e $\eta = 0{,}001$ não atingiu a
precisão em 200 épocas; assim, foi mantido $\eta = 0{,}01$.

\subsection{Validação e métricas}

Os modelos foram validados por simulações de Monte Carlo com $R = 10$
rodadas, conforme o item 5 do enunciado. Em cada rodada, as 640 amostras
foram embaralhadas e divididas em 80\% para treinamento (512 imagens) e
20\% para teste (128 imagens). A mesma partição foi usada pelos três
modelos em cada rodada, para que a comparação fosse justa, e a semente do
gerador de números aleatórios foi fixada para que os resultados possam ser
reproduzidos.

A medida de desempenho foi a acurácia no conjunto de teste, calculada a
partir da matriz de confusão implementada pela equipe: uma matriz
$C \times C$ em que o elemento $(i, j)$ conta as imagens da pessoa $i$
classificadas como pessoa $j$. A acurácia é a soma da diagonal dividida
pela soma de todos os elementos. Ao final das rodadas foram calculados a
média, o desvio-padrão, o maior e o menor valor da acurácia de cada modelo.
Para as rodadas de maior e menor acurácia foram plotadas as matrizes de
confusão (com a biblioteca Seaborn) e as curvas de aprendizado (EQM por
época). No Perceptron Simples, o EQM das curvas foi calculado com as saídas
do degrau bipolar.

\section{Reconhecimento Facial: Resultados e Discussão}\label{sec:res2}

\subsection{Desempenho geral}

A Tabela~\ref{tab:acc2} e a Fig.~\ref{fig:boxplot} mostram a acurácia dos
modelos nas 10 rodadas. A MLP obteve a maior média (97,42\%), seguida de
perto pelo Perceptron Simples (96,80\%). A diferença entre os dois (0,62
ponto percentual) é menor que o desvio-padrão de cada um (cerca de 1,1
ponto), e o maior e o menor valor foram iguais para ambos (0,9922 e
0,9531). Assim, nessas rodadas, os dois modelos tiveram desempenho muito
parecido. O ADALINE teve a menor média (92,50\%) e a maior variação entre
as rodadas, com acurácia entre 87,50\% e 96,09\%.

\begin{table}[htbp]
\caption{Estatísticas da acurácia ao longo das 10 rodadas}
\begin{center}
\begin{tabular}{|l|c|c|c|c|}
\hline
\textbf{Modelos} & \textbf{Média} & \textbf{Desvio-Padrão} & \textbf{Maior} & \textbf{Menor} \\
\hline
Perceptron Simples & 0,9680 & 0,0107 & 0,9922 & 0,9531 \\
\hline
ADALINE & 0,9250 & 0,0271 & 0,9609 & 0,8750 \\
\hline
MLP & 0,9742 & 0,0116 & 0,9922 & 0,9531 \\
\hline
\end{tabular}
\label{tab:acc2}
\end{center}
\end{table}

\begin{figure*}[htbp]
\centerline{\includegraphics[width=0.85\textwidth]{acuracia_boxplot.png}}
\caption{Boxplot (esquerda) e violin plot (direita) da acurácia de teste
nas 10 rodadas.}
\label{fig:boxplot}
\end{figure*}

O bom resultado de um modelo linear como o Perceptron Simples se explica
pela alta dimensão dos dados: há 900 atributos para apenas 512 imagens de
treinamento, e o Perceptron chegou a zero erros de treinamento em todas as
rodadas, ou seja, as 20 classes foram linearmente separáveis no conjunto de
treinamento. O ADALINE usa a mesma estrutura, mas minimiza o erro
quadrático da saída linear em vez de corrigir apenas as classificações
erradas, e não atingiu a precisão definida dentro do limite de épocas, o
que ajuda a explicar a sua menor acurácia. A MLP tem
$z = q(p+1) + m(q+1) = 50 \cdot 901 + 20 \cdot 51 = 46.070$ parâmetros, muito
mais que as 512 amostras de treinamento, o que vai contra a regra $N \gg z$
das notas de aula~\cite{b1} e indicaria risco de \textit{overfitting}. Mesmo assim, ela
teve a maior acurácia de teste, ou seja, generalizou bem.

\subsection{Matrizes de confusão}

As Figs.~\ref{fig:mc_ps}, \ref{fig:mc_adaline} e~\ref{fig:mc_mlp} mostram
as matrizes de confusão das rodadas de maior e menor acurácia. Como cada
pessoa tem cerca de seis imagens no conjunto de teste, os erros aparecem
como valores isolados fora da diagonal. Nas rodadas de maior acurácia, o
Perceptron Simples e a MLP erraram apenas uma imagem cada (phoebe
classificada como boland e mitchell classificada como danieln,
respectivamente), e o ADALINE errou cinco. Na rodada 7, a de menor acurácia
do Perceptron e da MLP, os dois modelos erraram seis imagens; cinco delas
eram das mesmas pessoas (an2i, glickman, karyadi, mitchell e saavik), e em
dois casos o erro foi idêntico (an2i classificada como sz24 e karyadi como
kawamura). Na rodada 8, a pior do ADALINE, os 16 erros ficaram espalhados
por 10 pessoas.

Somando as 10 rodadas, os erros se concentraram em poucas pessoas. A
pessoa mitchell foi a mais errada pela MLP (8 de 63 imagens de teste) e
pelo Perceptron (5 de 63), e glickman foi a classe mais escolhida por
engano por esses dois modelos. Por outro lado, ch4f nunca foi errada por
nenhum modelo, e a MLP acertou todas as imagens de 9 das 20 pessoas. O
ADALINE espalhou mais os erros: 18 das 20 pessoas tiveram ao menos uma
imagem errada.

Também houve um padrão nas imagens. Algumas foram erradas repetidamente
pelos três modelos, como a de kawamura olhando para cima e sorrindo (errada
quatro vezes por cada modelo) e a de choon virada para a esquerda, com
raiva e de óculos escuros (três vezes por cada modelo). Metade das imagens
do conjunto tem óculos escuros, mas eles aparecem em 64\% dos erros da MLP e
em 56\% dos erros do Perceptron. Na MLP, apenas 4 dos 33 erros foram em
rostos de frente. Ou seja, os erros desses dois modelos ocorrem mais quando
parte do rosto está coberta ou virada. No ADALINE isso não aconteceu com os
óculos (46\% dos erros), e os erros foram mais comuns nos rostos virados
para os lados (58 de 96).

\begin{figure*}[htbp]
\centerline{\includegraphics[width=0.9\textwidth]{confusao_perceptron.png}}
\caption{Matrizes de confusão do Perceptron Simples nas rodadas de maior
(esquerda) e menor (direita) acurácia. Linhas: classe real; colunas: classe
predita.}
\label{fig:mc_ps}
\end{figure*}

\begin{figure*}[htbp]
\centerline{\includegraphics[width=0.9\textwidth]{confusao_adaline.png}}
\caption{Matrizes de confusão do ADALINE nas rodadas de maior (esquerda) e
menor (direita) acurácia. Linhas: classe real; colunas: classe predita.}
\label{fig:mc_adaline}
\end{figure*}

\begin{figure*}[htbp]
\centerline{\includegraphics[width=0.9\textwidth]{confusao_mlp.png}}
\caption{Matrizes de confusão da MLP nas rodadas de maior (esquerda) e
menor (direita) acurácia. Linhas: classe real; colunas: classe predita.}
\label{fig:mc_mlp}
\end{figure*}

\subsection{Curvas de aprendizado}

A Fig.~\ref{fig:curvas} mostra as curvas de aprendizado das rodadas de
maior e menor acurácia. Nos três modelos, o EQM cai rapidamente nas
primeiras épocas e depois se estabiliza. O Perceptron Simples terminou o
treinamento sem erros entre 11 e 16 épocas em todas as rodadas. O ADALINE
sempre parou no limite de 100 épocas, com o EQM ainda acima de $10^{-3}$,
mas quase sem variação depois das primeiras dezenas de épocas. A MLP
atingiu a precisão entre 70 e 190 épocas, antes do limite de 200.

As curvas das rodadas de maior e menor acurácia são muito parecidas nos
três modelos. Logo, a diferença de acurácia entre as rodadas vem
principalmente de quais imagens caíram no conjunto de teste, e não de um
treinamento pior. Isso também aparece no fato de as rodadas de maior e
menor acurácia (6 e 7) serem as mesmas para o Perceptron Simples e para a
MLP.

\begin{figure*}[htbp]
\centerline{\includegraphics[width=\textwidth]{curvas_aprendizado.png}}
\caption{Curvas de aprendizado (EQM por época) nas rodadas de maior e menor
acurácia de cada modelo.}
\label{fig:curvas}
\end{figure*}

\section{Conclusão}\label{sec:conclusao}

\pendente{Conclusões da primeira etapa.}

Na segunda etapa, os três modelos reconheceram as 20 pessoas com acurácia
média acima de 92\%, mesmo com as imagens reduzidas para $30 \times 30$
pixels. A MLP teve a maior acurácia média (97,42\%), mas o Perceptron
Simples ficou muito próximo (96,80\%) e precisou de bem menos épocas, o que
mostra que, nesse problema de alta dimensão, um modelo linear já separa bem
as classes. O ADALINE teve o pior desempenho e a maior variação entre as
rodadas, além de não atingir a precisão definida dentro do limite de
épocas. As matrizes de confusão mostraram que os modelos tendem a errar as
mesmas pessoas e imagens, principalmente rostos com óculos escuros ou
virados, o que indica que essas imagens são mais difíceis de reconhecer
nessa resolução. Como trabalhos futuros, podem ser testadas as
outras dimensões de imagem, mais épocas para o ADALINE e um número maior de
rodadas de Monte Carlo.

\begin{thebibliography}{00}
\bibitem{b1} P. C. S. Barbosa, ``Notas de aula: Inteligência Artificial Computacional (T296),'' Universidade de Fortaleza, 2026, não publicado.
\bibitem{b2} F. Rosenblatt, ``The perceptron: a probabilistic model for information storage and organization in the brain,'' \textit{Psychological Review}, vol. 65, no. 6, pp. 386--408, 1958.
\bibitem{b3} B. Widrow and M. E. Hoff, ``Adaptive switching circuits,'' in \textit{IRE WESCON Convention Record}, part 4, pp. 96--104, 1960.
\bibitem{b4} D. E. Rumelhart, G. E. Hinton, and R. J. Williams, ``Learning representations by back-propagating errors,'' \textit{Nature}, vol. 323, no. 6088, pp. 533--536, 1986.
\bibitem{b5} T. M. Mitchell, \textit{Machine Learning}. New York: McGraw-Hill, 1997.
\end{thebibliography}

\end{document}
